\documentclass[11pt]{article}
\ifdefined\pdftrailerid\pdftrailerid{}\fi
\usepackage[T1]{fontenc}
\usepackage{lmodern}
\usepackage[margin=.88in]{geometry}
\usepackage{amsmath,amssymb,amsthm,mathtools,microtype}
\usepackage{xcolor}
\usepackage[colorlinks,linkcolor=blue!45!black,citecolor=blue!45!black,urlcolor=blue!55!black]{hyperref}
\hypersetup{pdftitle={Common Einstein cones and complete relational area probes for multifield matter},pdfsubject={Einstein--scalar matter, local characteristic cones and compact relational observables}}
\newtheorem{theorem}{Theorem}
\newtheorem{proposition}[theorem]{Proposition}
\newtheorem{lemma}[theorem]{Lemma}
\newcommand{\RR}{\mathbb R}
\newcommand{\dd}{\mathrm d}
\newcommand{\Lie}{\mathcal L}
\newcommand{\tr}{\operatorname{tr}}
\newcommand{\rank}{\operatorname{rank}}
\newcommand{\Id}{\mathrm{Id}}
\newcommand{\gm}{\mathfrak g}
\title{Common Einstein cones and complete relational\\area probes for multifield matter}
\author{}
\date{}
\begin{document}
\maketitle
\begin{abstract}
We classify smooth minimally coupled first-derivative scalar actions in four-dimensional Einstein gravity for which every physical scalar polarization propagates on the Einstein light cone on an attained open set. Strong hyperbolicity turns equality of characteristic sets into a matrix condition. Its gradient Hessian integrability forces a positive sigma model plus potential and, at rank four, at most a pullback target four-form. This applies to arbitrarily many scalar fields on the physical rank-four Gram manifold. A nonclosed four-form is invisible to the principal cone and direct stress, but can alter the field dynamics. We then show when ten types of marked spacelike area histories represent every compact coupled linear observable. Four complete material labels give an off-shell representative; additional scalars require reconstruction modulo the linearized equations. In an exact backreacting family with arbitrarily many scalars, this holds precisely when the normal-field mass matrix is cyclic from the metric readout. In a separate six-field model a stress-free derivative coupling makes an otherwise invisible scalar observable, without changing the background metric, direct stress response or principal cones.
\end{abstract}

\section{Introduction}
A shared characteristic cone is a restriction on the differential equations, not a guarantee that one can observe every physical field through a metric measurement. The distinction matters in Einstein gravity: an interaction can change scalar evolution without contributing either a principal symbol or a direct metric stress. We determine the most general such surplus term in a class of minimally coupled first-derivative actions, then give concrete tests for whether it helps or hinders a set of relational observables.

The earlier two-field cone theorem \cite{P46} isolates the Einstein metric cone, as opposed to a common effective scalar cone on a prescribed background. Examples of the latter appear in multifield DBI perturbation theory \cite{Langlois}. We use the Hessian in gradient variables to treat the full physical rank-four stratum, including \(N>4\) where the Gram variables lie on a proper rank manifold. Minors in first-order variational problems, and local Wess--Zumino forms, have established antecedents \cite{BCO,OS,DG}; the new restriction is that the open-data common-cone condition leaves precisely the specified four-form remainder. The result is local and does not require a choice of extension in the normal directions to the Gram manifold.

The observable part starts from a finite pointwise inverse for marked areas \cite{P47}, with actual material scalars defining the labels. This makes the linear observables relational in the sense of field-dependent coordinates \cite{Khavkine,BFR}. We prove completeness as a statement about compact coupled tests, not just a count of area types or a mode-separation argument. We use the term \emph{area history} for the whole set of compact test profiles of a fixed plane type over a spacetime label patch. Ten types do not mean ten numbers at a single time. The arbitrary-field cyclic mechanism has the usual linear-control algebraic structure; the derivative-only case uses the evolving geometry and an exact differential elimination, not a boundary observation inequality \cite{GKM}.

The genuinely curved data used to test the cone converse are given separately from the homogeneous examples used to test observability. All field equations below are classical, with area completeness at the linearized source or Peierls level on a regular patch. No global Hamiltonian stability, preferred quantum state, or nonlinear quantum reconstruction is asserted.

\section{Classification on the physical Gram stratum}\label{sec:cone}

We use signature \((-+++)\) and \(\kappa=8\pi G_N\). Let \(V_a^I=\partial_a\phi^I\), \(Q^{IJ}=-\tfrac12g^{ab}V_a^IV_b^J\). Consider a smooth first-derivative minimally coupled scalar density \(\sqrt{-g}L(\phi,V)\) with no curvature, extra spacetime tensor, or derivative-dependent change of frame. The parity-even class is \(L=P(\phi,Q)\); the oriented class may also use the volume form. The metric in these expressions is the metric of the Einstein term. Put \(q=\min(N,4)\). At the jets of interest \(V\) has rank \(q\). For \(N\le3\) assume independent gradient rows and a relatively open attained Gram image (the proof allows a degenerate induced Gram when its realizable neighborhood is among those tested). For \(N\ge4\) require an open attained family in the rank-four stratum of symmetric matrices with inertia \((1,3,N-4)\). Include \(\phi\) in these open families. In the oriented case use the corresponding oriented lift. On the tested family assume a positive normal scalar kinetic form, a strongly hyperbolic \(2N\)-companion for a metric-timelike normal, and precisely the Einstein metric cone as the real characteristic set.

\begin{theorem}[Common Einstein cone]\label{thm:cone}
Locally on the tested gradient-jet family, with a fixed orientation lift if needed, the matter density is
\begin{equation}\label{eq:form}
 \sqrt{-g}\bigl(G_{IJ}(\phi)Q^{IJ}-U(\phi)\bigr)\dd^4x+\phi^*A ,
 \qquad G=G^{\mathsf T}>0 ,
\end{equation}
where \(A\) is a target four-form. Conversely this form has the stated positive common principal cone. For \(N\le3\) the remainder is absent; for \(N=4\) it has no local bulk Euler contribution. For \(N>4\) the bulk force is governed by \(dA\), and it is Euler-null on an open full-rank family exactly when \(dA=0\) on the tested target patch. The direct stress of \(\phi^*A\) is zero in all cases.
\end{theorem}

\paragraph{From Gram data to an open calculation.}
At full row rank for \(N\le3\), equality of Gram makes the ordered rows define an isometry of their spans, including a possibly degenerate span. The finite-dimensional Witt extension theorem extends it to a Lorentz isometry on the ambient nondegenerate four-dimensional space. The Gram differential is onto all symmetric \(N\)-by-\(N\) matrices there: choose a right inverse of \(V g^{-1}\), then vary one factor in the symmetrized differential. At rank four choose an invertible four-row block \(V_B\). The same conclusion about the fibre follows by matching the Lorentz transform on \(V_B\); the remaining rows are fixed linear combinations with coefficients \(Q_{:B}Q_{BB}^{-1}\). Thus passing from relatively open invariant data to their corresponding tested jets gives open \(V\) boxes, using covariance; one restricts to the stated oriented lift where relevant. For \(N\le3\) it suffices also to work on the nondegenerate row spans and extend to any tested degenerate Gram by smoothness. We never differentiate \(P\) in fictitious normal directions to the deficient-rank stratum.

\begin{proof}[Proof of Theorem~\ref{thm:cone}]
Fix an orthonormal spacetime frame and \(\phi\), and write \(L_{Ia,Jb}\) for the gradient Hessian. The scalar principal pencil is \(-L_{Ia,Jb}k_a k_b\), with the sign such that \(GQ\) yields \(Gg^{ab}k_ak_b\). In harmonic gauge the Einstein block has the ordinary wave principal part; mixed metric--scalar variations here carry at most one differentiated variation and do not constrain the scalar polarizations at that order.

Normalize the timelike covector. For a unit spatial direction the only eigenvalues of the scalar companion are \(+1,-1\). Their multiplicities are constant on the connected sphere and are exchanged by reversing the spatial direction. They therefore both equal \(N\). Strong hyperbolicity implies semisimplicity, so each eigenspace has dimension \(N\); it is the kernel of the \(N\)-by-\(N\) pencil at the respective root. Hence the whole matrix vanishes there. By homogeneity it vanishes on the null cone. A real quadratic vanishing at \(k_0=\pm|\mathbf k|\) is proportional to the metric quadratic: subtraction kills the time-space cross terms, and addition fixes the Euclidean spatial block. Applying this entrywise gives
\begin{equation}\label{eq:hessian}
 \tfrac12\bigl(L_{Ia,Jb}+L_{Ib,Ja}\bigr)
       =-M_{IJ}(\phi,V)\eta^{ab},\qquad M>0.
\end{equation}
A repeated characteristic determinant alone would not give this matrix identity.

Let \(\epsilon_a=\eta^{aa}\) with no sum. The same-column Hessian is \(-\epsilon_a M_{IJ}\); the crossed spacetime indices of the off-diagonal Hessian in \eqref{eq:hessian} are antisymmetric. For any \(c\) choose \(a\ne c\). Commutation of third derivatives yields
\[
 \partial_{Kc}L_{Ia,Ja}
   =\partial_{Ia} L_{Kc,Ja}
   =-\partial_{Ia}L_{Ka,Jc}
   =-\partial_{Jc}L_{Ka,Ia}.
\]
Thus \(T_{IJK}=\partial_{Kc}M_{IJ}\) is symmetric in its first two and antisymmetric in its last two indices, which forces \(T=0\). This holds for every \(c\), so \(M=G(\phi)\).

Subtract \(GQ\). The remainder has zero same-column Hessians. It is separately affine in each gradient column on a local product box, hence a finite multiaffine polynomial. The crossed Hessian identities make the target coefficient antisymmetric in each pair of its selected column slots: compare coefficients of the other independent columns. Equivalently, every nonconstant coefficient multiplies a minor. At fixed target value, infinitesimal Lorentz invariance is a polynomial identity on this open box, hence an identity of polynomials. On spacetime exterior powers there are no Lorentz-invariant tensors in degrees one, two or three: spatial rotations already remove degree two, and boosts remove the time line and spatial volume left in degrees one and three. The degree-four invariant is the oriented volume form. The sole nonconstant remainder has the form \(A(\phi)(V_0,V_1,V_2,V_3)\); uniqueness under target-coordinate changes identifies \(A\) as a four-form. The constant is \(-U(\phi)\). Restoring the spacetime density gives \eqref{eq:form}.

The four-form has zero metric variation. Cartan's formula for a field variation \(W\) gives a bulk term \(\phi^*(\iota_WdA)\) and the exact boundary term \(\dd\phi^*(\iota_WA)\). The bulk term contains only first derivatives of the map, and it vanishes against a tangent variation \(W=V_a\) by antisymmetry. A five-form that evaluates to zero against all \(W\) and an open set of four independent columns is zero: the contraction is polynomial in those columns. For \(N=4\) the target five-form is identically zero, though boundary and global data may remain. Finally the sigma principal matrix is \(Gg^{-1}(k,k)\); positive \(G\) gives the usual independent scalar wave symmetrizer. The Einstein and gauge subsidiary principal parts are unchanged. This proves sufficiency without a nonlinear energy claim.
\end{proof}

\paragraph{Parity on rank four.}
On a rank-four block chart write \(M_{\mathrm{gr}}=Q_{:B}Q_{BB}^{-1}\). If \(A=\sum_{|S|=4}a_S(\phi)\dd\phi^S\) in increasing row sets, a local parity-even representative on that chart is
\begin{equation}\label{eq:gram}
 P=G_{IJ}Q^{IJ}-U+
 4\sqrt{-\det Q_{BB}}\sum_{|S|=4}a_S(\phi)\det(M_{\mathrm{gr},S}).
\end{equation}
Indeed \(V=M_{\mathrm{gr}}V_B\), and
\(\sqrt{-g}\sqrt{-\det Q_{BB}}=|\det V_B|/4\).
Thus \eqref{eq:gram} produces the pullback form or its negative on the two fixed sign lifts. In particular a locally parity-even Gram theory can have a nonclosed four-form force on a sign branch. Agreement of any two extensions from the physical Gram stratum gives identical pulled-back derivatives on the tested first jets; neither the principal operator nor the stress probes such an extension.

This branch statement differs from a reflection-reconnecting parity-even domain. At a fixed target value the polynomial coefficients just proved are unique on overlapping open boxes, so they propagate throughout a connected full-rank lift. If that lift also contains a \(V\) and its Lorentz reflection, the constant-coefficient four-form changes sign while \(Q\) and \(P\) do not. Equality for nearby pairs forces \(A=0\) on that component by the same open-set polynomial argument. The entire full-column-rank space is connected for \(N>4\) (deform to orthonormal columns and rotate one column using an unused target direction). It has two sign components for \(N=4\). This global obstruction does not prohibit the local representative \eqref{eq:gram}, and does not describe evolution through a lower-rank point.

\section{Attaining an open invariant set in the coupled equations}\label{sec:attain}
The classification concerns a tested open family, not merely an algebraic extension away from one solution. Here is a local constrained construction that attains the stated rank. It also gives a solution on which the extra term in Theorem~\ref{thm:cone} has a genuine bulk force.

Take \(G_{IJ}=\delta_{IJ}\), zero potential and cosmological constant. For \(N\ge5\) permit the oriented term
\begin{equation}\label{eq:attainA}
 A=\gamma\phi^5\dd\phi^1\wedge\dd\phi^2\wedge\dd\phi^3\wedge\dd\phi^4,
 \qquad \gamma\ne0 ;
\end{equation}
for \(N\le4\) the canonical model suffices. Thus the direct stress is exactly canonical. On a bounded spatial ball choose \(\chi=1\) near its centre and smooth with compact support in the ball, and prescribe
\[
 \phi^I=\phi_0^I+\epsilon\chi A_i^I x^i,\qquad
 \widetilde\Pi^I=\epsilon\chi u^I .
\]
Use conformally flat maximal data
\(\gamma_{ij}=\psi^4\delta_{ij}\),
\(K^{ij}=\psi^{-10}(\mathbb LW)^{ij}\),
\(\Pi^I=\psi^{-6}\widetilde\Pi^I\), where \(\mathbb L\) is the flat conformal Killing operator. With the convention for \(K\) in which the scalar flux enters with the displayed sign, the constraints are
\begin{align*}
 \partial_j(\mathbb LW)^{ij}
   &=\kappa\sum_I\widetilde\Pi^I\partial^i\phi^I,\\
 -8\Delta\psi
   &=\psi^{-7}\Bigl(|\mathbb LW|^2+\kappa\sum_I(\widetilde\Pi^I)^2\Bigr)
       +\kappa\psi\sum_I|\partial\phi^I|^2 .
\end{align*}
These are the standard conformal constraint weights \cite{CBIP}; the four-form contributes no constraint source. Impose \(W=0,\psi=1\) on the outer sphere as an auxiliary elliptic step. The vector Dirichlet operator has coercive energy on the ball, hence \(W=O(\epsilon^2)\). The scalar Dirichlet linearization at zero datum is \(-8\Delta\), and the implicit function theorem gives \(\psi=1+O(\epsilon^2)>0\), also in derivatives with respect to the finite parameters in a fixed bounded neighborhood. Interior evolution in harmonic gauge is the usual coupled quasilinear wave system with a smooth lower-order extra force. Its stress identity is consistent with the harmonic constraints since the extra force annihilates the target gradients. The auxiliary outer Dirichlet data do not prescribe timelike evolution boundary conditions; only a smaller interior domain is used.

At the centre put \(D=(u,A_i)\). In initial normal/spatial directions the realized metric Gram values are
\[
 \epsilon^{-2}Q=-\frac12D\,
       \operatorname{diag}(-\psi^{-12},\psi^{-4},\psi^{-4},\psi^{-4})D^{\mathsf T}.
\]
For \(N\ge4\), choose \(D\) near a frame with an invertible four-row block and sign fixed, for instance \(D=(I_4,0)^{\mathsf T}\). The stratum coordinates are the ten components of \(Q_{BB}\) and the \(4(N-4)\) components of \(Q_{\perp B}Q_{BB}^{-1}\). At \(\epsilon=0\) in the rescaled map, variations of the chosen invertible block span the symmetric entries, and independent variations of the other rows span its ratio coordinates. For \(N\le3\), the analogous full-row-rank frame has onto Gram differential by the argument of Section~\ref{sec:cone}. These submersions persist for fixed sufficiently small nonzero \(\epsilon\) by the parameter derivative estimate. The target values \(\phi_0^I\) can vary independently. This gives actual open attained data, on the physical rank manifold rather than in an ambient symmetric-matrix space.

In the five-field example take \(\phi^5=\Pi^5=0\) initially near the centre, with four independent remaining gradients and positive determinant choice. Normal second acceleration there is \(\gamma\det V_B/\sqrt{-g}\ne0\), since
\(\Box_g\phi^5+\gamma\det V_B/\sqrt{-g}=0\) in this orientation and its initial gradient vanishes. The nonclosed four-form therefore changes the solution even though its direct stress and principal symbol vanish. In a sign-fixed interior region the same result has the parity-even local description \eqref{eq:gram}. This is an attainment benchmark for explicit actions, not a claim that every allowed action attains all invariant data.

\section{Compact source completeness from marked area histories}\label{sec:area}
Take a regular material chart \(X\) on a solution, supplied by natural local functions of the fields, and let \(\gm=(X^{-1})^*g\) on a chosen inverse branch. Its three spatial coordinate directions are assumed spacelike with a positive margin on a shrunk measurement patch. For constant labelled tangents \(v,w\) define
\[
 a(v,w)=\sqrt{\gm(v,v)\gm(w,w)-\gm(v,w)^2},\qquad
 A_I(f)=\int f(X)a_I(X)\,\dd^4X .
\]
Only spacelike pairs are used. This is a compactly smeared observable along a family of parallel labelled two-surfaces, with no unsmeared surface operator involved.

\begin{lemma}[Ten plane types]\label{lem:area}
The following ten area types have a smooth pointwise inverse for \(\gm\) sufficiently locally, with a fixed small nonzero tilt on the reference patch. Their differential has a zero-order left inverse on symmetric metric perturbations.
\end{lemma}
\begin{proof}
In the three spatial labels let \(q=(\gm_{ij})>0\). Under the coordinate alternating-symbol map from bivectors to three components \(n\), the squared area is \(n^{\mathsf T}C n\), \(C=(\det q)q^{-1}\). The six \(n=e_i,\ e_i+e_j\) determine \(C\) by polarization, hence \(q=\sqrt{\det C}C^{-1}\). Write \(s_i=\gm_{0i}\), \(v=\gm_{00}\). For cyclic \((i,j)=(1,2),(2,3),(3,1)\) tilt the plane to \((e_i\pm\alpha e_0,e_j)\). Its squared area is
\[
 L_i^\pm=L_i^0\pm2\alpha(q_{jj}s_i-q_{ij}s_j)
              +\alpha^2(q_{jj}v-s_j^2),\quad L_i^0=q_{ii}q_{jj}-q_{ij}^2.
\]
Keep only \(L_1^\pm,L_2^+,L_3^+\). The pair gives \(d_1=q_{22}v-s_2^2\), so \(v=(d_1+s_2^2)/q_{22}\), as well as the first signed shift equation. Divide the other two one-sided differences by \(2\alpha\). The Jacobian in the \(s\)'s is the cyclic linear matrix \((q_{jj}s_i-q_{ij}s_j)_{i=1}^3\), plus \(O(\alpha)\) on the branch. Its determinant is \(q_{11}q_{22}q_{33}-q_{12}q_{23}q_{31}>0\), by the positive two-dimensional principal minors. The implicit function theorem solves the exact shifted equations for fixed small \(\alpha\ne0\). Spacelikeness holds by the strict margin and continuity. Differentiating the exact inverse gives the claimed zero-order inverse; this is the ten-plane specialization of \cite{P47}.
\end{proof}

Let \(u=(h,\pi)\) be a coupled perturbation and \(D\xi=(\Lie_\xi g,d\phi\,\xi)\). The variation at fixed material labels is
\begin{equation}\label{eq:projection}
 \widehat u=u-DZ,\qquad Z=(dX)^{-1}X'[u].
\end{equation}
Naturality gives \(X'[D\xi]=dX\,\xi\), so \eqref{eq:projection} is a gauge-invariant projection, expressed on the reference chart. A compact coupled source \(t=(t_g,t_\phi)\) is invariant precisely in the linear Ward sense \(D^*t=0\). All adjoints below include the appropriate density relative to the fixed-label pairing.

\begin{theorem}[Off-shell and on-shell source conversion]\label{thm:source}
On a regular chart as above:
\begin{enumerate}
\item If all material fields can be read as \(\phi=F(X)\) on nearby fields for a fixed function \(F\), every compact Ward source has exactly the same \emph{off-shell} pairing as compact profiles of the ten area types.
\item Let \(J\) be the coupled Jacobi operator. More generally, if the extra fixed-label variations obey
\begin{equation}\label{eq:identity}
 \widehat\pi=\mathcal K\widehat h+\mathcal L J u
\end{equation}
as a local differential identity, then the same ten types give every compact coupled source class \emph{modulo the linearized equations}. The profiles do not enlarge support.
\end{enumerate}
These statements pass to the local physical linear algebra under the standard causal, hyperbolic-gauge and quotient hypotheses on a causally convex patch or a stipulated ambient problem.
\end{theorem}
\begin{proof}
One has \(\langle t,u\rangle=\langle t,\widehat u\rangle\) by \(D^*t=0\); a compact cutoff on \(Z\) near \(\operatorname{supp}t\) suffices to justify the pairing. In (1) differentiate the off-shell material identity at fixed labels to get \(\widehat\pi=0\). In (2) instead obtain
\[
 \langle t,u\rangle
 =\langle t_g+\mathcal K^*t_\phi,\widehat h\rangle
     +\langle\mathcal L^*t_\phi,Ju\rangle ,
\]
where the already-fixed label fields have no variations in this expression. Formal adjoints of local operators preserve compact support. If \(M_{\rm area}\) is the map from \(\widehat h\) to the ten direct area variations and \(L_{\rm area}M_{\rm area}=1\) is the differential inverse from Lemma~\ref{lem:area}, use
\(f=L_{\rm area}^*(t_g+\mathcal K^*t_\phi)\) on the label patch, with the Jacobian density from the change of coordinates. This gives the stated pairings. Conversely any compact area profile is Ward closed by \eqref{eq:projection}, and the second term pairs trivially in the equation quotient. Equivalence of the gauge-fixed wave formulation is used to define the causal algebra, not inferred from this material projection.
\end{proof}

For \(X=\phi\) a regular quartet, with background volume density understood, the Ward identity is
\(-2\nabla_\mu t_g{}^\mu{}_\nu+t_I\partial_\nu\phi^I=0\).
Its scalar component is determined by the arbitrary compact \(t_g\) since \(d\phi\) is invertible. This is the kinematic meaning of (1); it does not delete the scalar contribution to the field equations.

There is a useful direct criterion for (2). In Gaussian-normal target coordinates around an embedded material background \(Y=0\), use its tangential coordinates as the actual four labels. Let \(K_{A\mu\nu}\) be the target second fundamental tensor, with \(\partial_A G_{\mu\nu}=-2K_{A\mu\nu}\) on the surface. At fixed labels the trace-reversed Einstein residual linearizes to
\begin{equation}\label{eq:direct}
 \delta{\cal E}_{\mu\nu}
   =\kappa^{-1}\delta R_{\mu\nu}[\widehat h]
     -U(X,0)\widehat h_{\mu\nu}-C_{A\mu\nu}\pi^A,\qquad
 C_{A\mu\nu}=-2K_{A\mu\nu}+(\partial_AU)g_{\mu\nu}.
\end{equation}
Here \(G_{\mu A}(X,0)=0\) and the background normal gradients vanish. The four-form has no direct stress term. An injective \(C\) on the patch supplies (2) by a smooth algebraic left inverse, without inverting a wave operator. At most ten normal fields can be recovered this way directly; there is no such bound for the dynamical reconstruction below.

\section{Exact backreacting matter and hidden-field reconstruction}\label{sec:examples}

\subsection{A material FLRW branch with an active four-form}
For the canonical massless quartet, the exact expanding flat FLRW solution satisfies
\begin{equation}\label{eq:FLRW}
 \begin{gathered}
 \dd s^2=-\dd t^2+a(t)^2\dd\mathbf x^2,\qquad
 \dot T=P a^{-3},\qquad X^i=\ell x^i,\qquad P,\ell>0,\\
 \rho=P^2/(2a^6)+3\ell^2/(2a^2),\qquad
 \dot a/a=\sqrt{\kappa\rho/3},\qquad a(0)=1.
 \end{gathered}
\end{equation}
The four waves and stress conservation imply the remaining Einstein equation. The matter density is nonzero and \(X=(T,X^i)\) is a regular material chart. For example in the uncoupled quartet varying \(P\) at \(a(0)=1,T(0)=0\) changes the marked inverse time component \(-P^2\) at the same labelled event. This gives an observable metric change in exact solutions, rather than a passive recalibration of the clock. For the fixed actions below we use a chosen solution \eqref{eq:FLRW}.

Add \(m\ge1\) extra scalars \(Y^A\), with metric, potential and oriented four-form
\begin{align}
 G&=(1+2bY^1)(\dd T)^2+\sum_{i=1}^3(\dd X^i)^2
                         +\sum_A(\dd Y^A)^2,\quad b\ne0,\qquad
 U=\tfrac12 Y^{\mathsf T}M Y, \label{eq:target}\\
 A_0&=\gamma Y^1\,\dd T\wedge\dd X^1\wedge\dd X^2\wedge\dd X^3,\qquad
 \gamma=-bP/\ell^3 . \label{eq:WZ}
\end{align}
Here \(M\) is constant real symmetric, which may be taken positive, and \(G>0\) on a small target tube. The coordinate labels themselves are dynamical scalars. The oriented form has the local sign-fixed Gram description \eqref{eq:gram}. In material coordinates on \eqref{eq:FLRW},
\[
 \bar g=-(a^6/P^2)(\dd T)^2+(a^2/\ell^2)\sum_i(\dd X^i)^2,\qquad
 \sqrt{-\bar g}\,\bar g^{00}=-P/\ell^3 .
\]
Writing \(j_X=\det(\partial X)/\sqrt{-g}\) in a positively oriented coordinate chart, the first normal equation is
\[
 E_1=\Box_gY^1-(MY)_1-b\,g^{-1}(\dd T,\dd T)+\gamma j_X=0 .
\]
The last two terms cancel on the chosen \(Y=0\) FLRW branch by \eqref{eq:WZ}. All other normal equations hold there. The tangential scalar variations of the added terms vanish there since \(Y=\dd Y=0\). The stress is exactly the canonical quartet's. Thus \eqref{eq:FLRW} with \(Y=0\) is an exact coupled solution for \eqref{eq:target}--\eqref{eq:WZ}, not a fixed-background solution with omitted normal equations. Here \(dA_0\ne0\), and the full gradient rank is four. The characteristic cone is the Einstein one, by Theorem~\ref{thm:cone} or directly by the positive sigma principal form.

\subsection{Arbitrarily many scalars: a necessary and sufficient criterion}
On this solution use fixed labels and write the metric and normal perturbations as \(h,\pi\). Let \(e=\delta{\cal E}_{00}\) be the indicated trace-reversed Einstein residual of \eqref{eq:direct}, and let \(r=\delta E\) be the vector of normal residuals. At \(Y=0\) the fixed-label operator is exactly
\begin{align}
 e&=\kappa^{-1}\delta R_{00}[h]-2b\pi_1,\qquad
 r=(\bar\Box\Id-M)\pi+e_1 F[h], \label{eq:resclass}\\
 F[h]&=b\bar g^{0\mu}\bar g^{0\nu}h_{\mu\nu}
                   -\tfrac12\gamma\bar j\,\tr_{\bar g}h,\qquad
 \bar j=1/\sqrt{-\bar g}. \nonumber
\end{align}
The first unit column \(e_1\) refers to the normal target; \(e\) denotes the residual. There is no metric variation of \(\Box_gY\) on the constant \(Y\) background. The \(F\) term comes from \(\delta g^{00}=-\bar g^{0\mu}\bar g^{0\nu}h_{\mu\nu}\) and \(\delta j_X=-\bar j\tr_{\bar g}h/2\); the potential has no first-order metric stress variation. Only \(\pi_1\) is directly read by the metric.

\begin{theorem}[Cyclic completeness and its obstruction]\label{thm:cyclic}
For \eqref{eq:target}--\eqref{eq:resclass}, the ten area histories represent all compact coupled linear source classes if and only if \(e_1\) is cyclic for \(M\), i.e. \(\operatorname{span}\{M^je_1:j\ge0\}=\RR^m\). This occurs for choices of positive \(M\) at every finite \(m\). More generally the cyclic normal subspace is reconstructible, and its orthogonal complement supports invisible Jacobi waves.
\end{theorem}
\begin{proof}
Put \(c_j=e_1^{\mathsf T}M^j\pi\). Once \(c_0=\pi_1\) is read from \(e,h\), \eqref{eq:resclass} gives
\[
 c_{j+1}=\bar\Box c_j-e_1^{\mathsf T}M^j r
                             +(e_1^{\mathsf T}M^je_1)F[h].
\]
Here \(M\) is constant so commutes with the common wave operator. Invert the finite row matrix on the cyclic subspace using its first \(\dim\) independent rows. This is exactly a support-local differential formula of the type \eqref{eq:identity}, with no inverse wave operator and order at most twice the cyclic dimension in \(h\). A positive symmetric tridiagonal matrix with every adjacent off-diagonal nonzero and a sufficiently large diagonal realizes full cyclicity: \(M^je_1\) has no components later than \(j+1\), where its entry is the product of the \(j\) adjacent off-diagonals.

Conversely, the symmetry of \(M\) makes the orthogonal complement invariant. Let \(w\ne0\) be an eigenvector there, \(Mw=\lambda_w w\), and choose a nonzero smooth local solution \(f\) of \((\bar\Box-\lambda_w)f=0\). The perturbation \(h=\delta X=0,\ \pi=wf\) satisfies the full coupled Jacobi equations. Indeed \(w_1=0\), so the metric and tangential variations are unforced, including the variation of \eqref{eq:WZ}, which is through the first normal and its gradient. The normal residual vanishes by \eqref{eq:resclass}. Since \(dY=0\) on the background this mode is not diffeomorphism gauge. A compact normal test along \(w\) is Ward closed and can pair nontrivially, whereas every marked area history vanishes on this mode. It therefore cannot be represented by an area profile even modulo the Jacobi equations. In the causal construction, take \(f\) to be the Klein--Gordon causal propagator applied to a compact scalar test with nonzero image. The invisible eigendirection is a decoupled scalar wave sector, so a second compact Ward test detects that image by a nonzero physical Peierls pairing. The obstruction is not removed by quotienting the presymplectic radical.
\end{proof}

For clarity the single-output assumption has a constant-matrix extension within the same exact FLRW construction. Replace the tangential target block by \(\delta_{\mu\nu}+2\sum_A Y^A B_{A\mu\nu}\) with constant symmetric \(B_A\); take the coefficient of \(Y^A\dd^4X\) in \eqref{eq:WZ} to be
\(\gamma_A(T)=\sqrt{-\bar g(T)}B_{A\mu\nu}\bar g^{\mu\nu}(T)\).
These are fixed target functions of the dynamical clock, not external tensors; shrink the tube so \(G>0\). The same FLRW quartet then solves the normal balance, and the direct linear readout is \(C\pi=2\sum_A B_A\pi^A\). Successive local observables \(CM^j\pi\) replace \(c_j\), so the criterion is \(\operatorname{span}_j M^j\operatorname{im}C^{\mathsf T}=\RR^m\) by the same proof. A hidden eigenvector annihilating \(C\) also annihilates \(\gamma_A(T)\) in its combination for every \(T\), so it supplies the same invisible witness. This realizes the direct-injective regime for \(m\le10\) and lower-rank but dynamically complete regimes for larger \(m\). It does not generalize the criterion to arbitrary field-dependent or noncommuting operators.

\subsection{The derivative four-form alone can enable completeness}
A separate six-field construction (\(m=2\)) makes the observation mechanism depend on the direct-stress-free term itself. Use diagonal masses \(M=\operatorname{diag}(\mu_1^2,\mu_2^2)\) with \(\mu_2^2>0\), and add to \eqref{eq:WZ}
\begin{equation}\label{eq:gyroscopic}
 A_1=\frac{\beta}{2}(Y^1\dd Y^2-Y^2\dd Y^1)\wedge
                      \dd X^1\wedge\dd X^2\wedge\dd X^3 .
\end{equation}
The background force of \(dA_1\) vanishes because it needs a normal gradient; the first-order tangential force also vanishes there. The exact background, direct metric stress response, and principal cones are identical for all \(\beta\). With fixed labels the change in the linear normal equations is
\[
 \begin{split}
 r_1&=(\bar\Box-\mu_1^2)\pi_1+F[h]+\beta\bar j\,\partial_T\pi_2,\\
 r_2&=(\bar\Box-\mu_2^2)\pi_2-\beta\bar j\,\partial_T\pi_1 ,
 \end{split}
\]
as follows by varying the oriented density
\(\tfrac\beta2(Y^1\partial_TY^2-Y^2\partial_TY^1)\).
At \(\beta=0\) the second normal field is invisible by Theorem~\ref{thm:cyclic}. At \(\beta\ne0\), resolve from these equations and the Einstein readout the two local expressions
\[
 d=\frac{r_1-(\bar\Box-\mu_1^2)\pi_1-F[h]}{\beta\bar j}
       =\partial_T\pi_2,\qquad
 q=r_2+\beta\bar j\,\partial_T\pi_1=(\bar\Box-\mu_2^2)\pi_2 .
\]
The clock \(T\) is harmonic, so
\(\bar\Box=-P^2 a^{-6}\partial_T^2+\ell^2a^{-2}\Delta_X\)
has no time first derivative. Put \(B=a^2/\ell^2\), \(C=P^2/(a^4\ell^2)\). Then \(Bq=-C\partial_T d+\Delta_X\pi_2-\mu_2^2 B\pi_2\). Differentiating once in \(T\) eliminates the unknown spatial Laplacian because \(\partial_T\pi_2=d\):
\begin{equation}\label{eq:derive}
 \pi_2=\frac{-\partial_T(C\,\partial_Td)+\Delta_Xd
                   -\mu_2^2Bd-\partial_T(Bq)}
                 {\mu_2^2\partial_TB}.
\end{equation}
The denominator is nonzero on the expanding branch of \eqref{eq:FLRW}. Equation \eqref{eq:derive} is an identity modulo local residuals of the form \eqref{eq:identity}; no initial condition, boundary inversion or time integration enters it. This proves completeness for \(\beta\ne0\) with the stated mass and branch. In contrast if \(\mu_2=0\), a constant pure second normal variation is still invisible even at \(\beta\ne0\). In this example it is the combination of the derivative four-form, mass and expansion that enables the readout, rather than a generic property of metric-insensitive couplings.

\section{A curved constraint and a conformally transported scalar}\label{sec:curved}
This section is a separate illustration of the source criterion, not an extension of the minimal classification to nonminimal theories or a curved version of the flat spectral result in \cite{P45}.

For a reconstructed marked metric write \(q_{ij}=\gm_{ij}\), \(\beta^i=q^{ij}\gm_{0j}\), \(N=(q^{ij}\gm_{0i}\gm_{0j}-\gm_{00})^{1/2}\), and \(K_{ij}=(\partial_tq_{ij}-(\Lie_\beta q)_{ij})/(2N)\). The time label may be fixed by a material recalibration to agree with background proper time; it is not then replaced by an external coordinate choice in perturbed solutions. Gauss--Codazzi gives the local geometric normal constraint
\begin{equation}\label{eq:hc}
 {\cal C}_g=\kappa^{-1}G_g(n,n)
   =\frac{{}^{(3)}R[q]+K^2-K_{ij}K^{ij}}{2\kappa}.
\end{equation}
The full material chart also fixes the true normal sigma energy
\[
 \rho_\phi=\tfrac12G_{IJ}\bigl(\Pi^I\Pi^J+q^{ij}\partial_i\phi^I\partial_j\phi^J\bigr)
                 +U,\qquad
 \Pi^I=N^{-1}(\partial_t\phi^I-\beta^i\partial_i\phi^I),
\]
and the four-form contributes no term. Compact variations of these local quantities and of their difference are area-history tests by support-preserving adjoints. Their difference vanishes in the constrained theory. Instantaneous areas alone need not give the off-shell geometric functional: on the flat-slice FLRW reference with \(H=\dot a/a\ne0\) take \(h_{ij}=2a^2f(t)\delta_{ij},h_{0\mu}=0\), \(f(t_0)=0\), \(\dot f(t_0)\ne0\). Instantaneous areas at \(t_0\) do not vary, but \(\delta K=3\dot f(t_0)\) and \(\delta G(n,n)=2H\,\delta K\ne0\). This is not an obstruction to instantaneous constrained energy readout.

For a different model take an Einstein-frame metric \(\widetilde g\), the canonical quartet of \eqref{eq:FLRW}, and an additional scalar \(\sigma\) with Lagrangian
\(-\tfrac12(\widetilde g^{-1}(\dd\sigma,\dd\sigma)+\tfrac16\widetilde R\,\sigma^2)\).
The exact branch has \(\sigma=0\), so the conformal scalar stress has no linear part there. Measure instead the Jordan metric \(g_J=e^{\kappa f\sigma}\widetilde g\), \(f\ne0\). The same four rods are transported with the action under this change of variables; they are not independently chosen to be Jordan-minimal. Use fixed material labels and \(O=f\,\delta\sigma\). Then \(\widehat h_E=\widehat h_J-\kappa O g_0\), where \(g_0\) is the common background metric. Let \(J_E\) be the linearization of \(G-\kappa T^\phi\) at fixed rods. The covariant massless rod stress at fixed rods is unchanged under this Weyl variation; \(g_{ab}g^{cd}\sum_I\partial_c\phi^I\partial_d\phi^I\) is unchanged. Thus
\[
 J_E(\widehat h_J)=J_E(\widehat h_E)
                     +\kappa(g_0\Box-\nabla\nabla)O.
\]
On shell \(J_E(\widehat h_E)=0\) and \((\Box-\mathcal R/6)O=0\), with \(\mathcal R=R[g_0]\). Put
\[
 \begin{split}
 W_{\mu\nu}&=J_E(\widehat h_J)_{\mu\nu},\qquad
 \tau=2\tr_{g_0} W/\kappa,\\
 B_\mu&=-\kappa^{-1}g_{0,\mu\lambda}(R^{-1})^{\lambda\nu}
                                      \nabla^\alpha W_{\alpha\nu}.
 \end{split}
\]
On the active quartet solution \(R_{\mu\nu}=\kappa\sum_I\partial_\mu\phi^I\partial_\nu\phi^I\) is a positive-definite covariant form and therefore invertible, even where its Lorentzian trace vanishes. Covariant differentiation yields \(\nabla^\mu W_{\mu\nu}=-\kappa R_{\nu\lambda}\nabla^\lambda O\) on shell, and the trace gives \(\tau=\mathcal R O\). Thus \(B=\dd O\), \(\dd\tau-\mathcal RB=O\dd\mathcal R\). For a smooth positive contravariant form \(b\) and fixed positive length scale \(L_0\), locally reconstruct
\begin{equation}\label{eq:curv}
 O=\frac{\mathcal R\,\tau+
      L_0^2 b(\dd\mathcal R,\dd\tau-\mathcal RB)}
          {\mathcal R^2+L_0^2 b(\dd\mathcal R,\dd\mathcal R)} .
\end{equation}
The denominator need not exclude curvature zeros: on \eqref{eq:FLRW} one has
\(\mathcal R=\kappa(3\ell^2/a^2-P^2/a^6)\);
at a zero its derivative in background proper time is \(12\kappa H\ell^2/a^2\ne0\). Off shell the difference \(\tau-\mathcal R O\) is
\(2\tr_{g_0}J_E(\widehat h_E)/\kappa+6(\Box-\mathcal R/6)O\), and \(B-\dd O\) is the inverse-Ricci contraction displayed above with \(W\) replaced by \(J_E(\widehat h_E)\). Inserting these in \eqref{eq:curv} shows explicitly that the error is a local coupled Jacobi combination. Theorem~\ref{thm:source} therefore applies to the fifth scalar as well, on shell, through the simple curvature crossing.

On the scalar mode \(\widehat h_E=0\), the Jordan geometric readout in \eqref{eq:hc} is \(\delta{\cal C}_J=3H\dot O-a^{-2}\Delta O-3H^2O\); the last term is from normalizing the Jordan unit normal. The Einstein-frame geometric readout on this mode is zero, as is the full constrained linear combination. This case distinguishes a nonzero bare geometric response from a complete constraint. It is local, in a specified transported frame, and is not an assertion about a Poincar\'e vacuum atom, a charge at infinity, or an arbitrary scalar-flat region.


\begin{thebibliography}{15}\setlength{\itemsep}{2pt}
\bibitem{P45} \emph{A massless stress pole from the variance of a spatially scaled vacuum energy}, research manuscript 45, \href{https://github.com/g3173028019-beep/ai-automated-quantum-gravity-research/tree/main/papers/45-vacuum-stress-massless-pole}{source and manuscript}.
\bibitem{P46} \emph{When Two Scalar Modes Share the Einstein Light Cone}, research manuscript 46, \href{https://github.com/g3173028019-beep/ai-automated-quantum-gravity-research/tree/main/papers/46-einstein-scalar-common-cone}{source and manuscript}.
\bibitem{P47} \emph{Spacelike area tomography and compact linear gravitational probes}, research manuscript 47, \href{https://github.com/g3173028019-beep/ai-automated-quantum-gravity-research/tree/main/papers/47-spacelike-area-tomography}{source and manuscript}.
\bibitem{Langlois} D.~Langlois, S.~Renaux-Petel, D.~A.~Steer and T.~Tanaka, \emph{Primordial perturbations and non-Gaussianities in DBI and general multi-field inflation}, Phys. Rev. D \textbf{78} (2008), 063523, \href{https://arxiv.org/abs/0806.0336}{arXiv:0806.0336}.
\bibitem{BCO} J.~M.~Ball, J.~C.~Currie and P.~J.~Olver, \emph{Null Lagrangians, weak continuity, and variational problems of arbitrary order}, J. Funct. Anal. \textbf{41} (1981), 135--174, \href{https://doi.org/10.1016/0022-1236(81)90085-9}{publisher}.
\bibitem{OS} P.~J.~Olver and J.~Sivaloganathan, \emph{The structure of null Lagrangians}, Nonlinearity \textbf{1} (1988), 389--398.
\bibitem{DG} J.~Davighi and B.~Gripaios, \emph{Homological classification of topological terms in sigma models on homogeneous spaces}, JHEP \textbf{09} (2018), 155, \href{https://arxiv.org/abs/1803.07585}{arXiv:1803.07585}.
\bibitem{CBIP} Y.~Choquet-Bruhat, J.~Isenberg and D.~Pollack, \emph{The Einstein--scalar field constraints on asymptotically Euclidean manifolds}, Chin. Ann. Math. Ser. B \textbf{27} (2006), 31--52, \href{https://arxiv.org/abs/gr-qc/0506101}{arXiv:gr-qc/0506101}.
\bibitem{Khavkine} I.~Khavkine, \emph{Local and gauge invariant observables in gravity}, Class. Quantum Grav. \textbf{32} (2015), 185019, \href{https://arxiv.org/abs/1503.03754}{arXiv:1503.03754}.
\bibitem{BFR} R.~Brunetti, K.~Fredenhagen and K.~Rejzner, \emph{Quantum gravity from the point of view of locally covariant quantum field theory}, Commun. Math. Phys. \textbf{345} (2016), 741--779, \href{https://arxiv.org/abs/1306.1058}{arXiv:1306.1058}.
\bibitem{GKM} S.~Gerbi, C.~Kassem, A.~Mortada and A.~Wehbe, \emph{Exact controllability and stabilization of locally coupled wave equations: theoretical results}, \href{https://arxiv.org/abs/2003.14001}{arXiv:2003.14001}.
\end{thebibliography}
\end{document}
